Building on the full-montage reference in Section~\ref{sec:exp1_reference}, the complete Stage~A--Stage~C Proposed-Med pipeline was rerun independently on each channel subset in both Internal and Cao2018, using 30-second epochs and the median centre. Internal comprised 46 sessions and 20 channel subsets, including 14 anatomical subsets (the full montage among them) and 6 single-electrode subsets, whereas Cao2018 comprised 58 sessions and 18 channel subsets, including 14 anatomical subsets (the full montage among them) and 4 single-electrode subsets; in each corpus, 13 anatomical subsets were compared against the full montage. In every subset rerun, Stage~A, Stage~B and Stage~C all operated only on the electrodes available within that subset, so each subset constituted an independent, self-contained detector rather than the full-montage pipeline with channels subsequently hidden.

Because Stage~C continued to score every electrode retained in a subset individually, each subset rerun yielded one $F_1$ value per electrode per session rather than a single subset-level score. This per-electrode scoring makes it possible to examine detection performance at the level of individual scalp locations before those results are aggregated into the region-mean statistics reported later in this section. Internal electrodes are reported below using their corresponding 10--20 scalp locations.

Within these self-contained subset reruns, per-electrode detection quality remained concentrated at the frontopolar and neighbouring anterior frontal electrodes rather than distributed uniformly across the frontal region (Figure~\ref{fig:exp1_subset_per_electrode}). On Internal, the four highest-performing electrodes across the anatomical subset reruns were Fp2 ($F_1=85.28\%$), Fp1 ($84.84\%$), AF4 ($77.04\%$) and AF3 ($70.70\%$), all scored within the frontal subset's own rerun; performance then fell sharply at the two remaining frontal electrodes, F3 and F4, to $42.68\%$ and $41.00\%$, respectively. On Cao2018, the four highest-performing electrodes were Fp1 ($78.80\%$), Fp2 ($77.05\%$), F3 ($55.97\%$) and F4 ($54.64\%$), which together made up the entire Cao2018 frontal subset, with the frontopolar pair leading F3 and F4 by roughly 22 percentage points. Although every one of these frontal electrodes scored lower in its own subset rerun than the same electrode did in the full-montage run (Section~\ref{sec:exp1_reference}), the ranking of electrodes by $F_1$ was unchanged from the full-montage run in both datasets, so restricting Stage~A/B screening to the frontal subset alone did not reorder the strongest frontal electrodes.

The central electrodes again showed a dataset-specific asymmetry within their own subset rerun. On Internal, C3 reached $F_1=58.59\%$ while C4 reached only $F_1=1.73\%$, so the pronounced C3--C4 asymmetry seen in the full-montage run persisted in the central subset's self-contained rerun; on Cao2018, C3 and C4 remained close together, at $27.67\%$ and $26.62\%$, respectively, so this asymmetry stayed specific to Internal. Parietal and occipital electrodes were much lower in their own subset reruns: the best-performing parietal electrode was CP1 at $1.92\%$ on Internal and CP4 at $15.63\%$ on Cao2018, and the best-performing occipital electrode was PO3 at $3.35\%$ on Internal and P3 at $10.04\%$ on Cao2018. Overall, the self-contained subset reruns preserved the same anterior-to-posterior spatial pattern observed under the full montage, with detection concentrated at the frontopolar and anterior frontal sites and much lower performance at the central, parietal and occipital electrodes.

\begin{figure}[ht]
    \centering
    \includegraphics[width=0.9\linewidth]{fig_exp1_subset_per_electrode.pdf}
    \caption{Per-electrode detection $F_1$ within each region's own channel-subset rerun,
        coloured by scalp region. Each bar is one electrode scored inside that region's
        own self-contained subset run (Table~\ref{tab:exp1_subset_summary}), not inside
        the 32-channel run of Figure~\ref{fig:region_performance} and not a
        single-electrode pipeline.}
    \label{fig:exp1_subset_per_electrode}
\end{figure}

The descriptive differences between the subset reruns and the full-montage reference were next tested at the individual-electrode level using session-paired comparisons. For every electrode entering one of the four anatomical subsets shown in Figure~\ref{fig:exp1_subset_per_electrode}, that electrode's session-level $F_1$ inside its own subset rerun was paired against the session-level $F_1$ of the same electrode inside the full-montage run, on the sessions common to both runs; every electrode retained its dataset's full session pool (46 for Internal, 58 for Cao2018). Each electrode was tested with a two-tailed Wilcoxon signed-rank test, paired by session, with $p$ Bonferroni-corrected over every electrode tested within each dataset (24 for Internal, 18 for Cao2018, a family separate from the 13 region-mean tests below) and the matched-pairs rank-biserial correlation reported as the effect size, signed so that a negative value again indicates lower $F_1$ in the subset rerun (Table~\ref{tab:exp1_subset_per_electrode_stats}).

On Internal, 11 of the 24 electrodes tested remained significantly different after correction, and every significant difference was negative. Despite the visible drop in the descriptive means, no statistically significant difference was detected for either of the two highest-performing electrodes: Fp2 changed from $F_1=86.82\%$ in the full montage to $F_1=85.28\%$ in the frontal subset rerun (adjusted $p=1.000$, $r_{\mathrm{rb}}=-0.05$), and Fp1 changed from $86.32\%$ to $84.84\%$ (adjusted $p=1.000$, $r_{\mathrm{rb}}=-0.08$). The remaining two of the four highest-performing electrodes, AF4 and AF3, did reach significance, with reductions that were statistically significant but modest in absolute magnitude, $-3.28$ and $-3.06$ percentage points, respectively (adjusted $p=0.003$ and $p<0.001$; $r_{\mathrm{rb}}=-0.67$ and $-0.69$). Among the non-frontal electrodes, C3 showed a significant reduction ($-2.47$ points, adjusted $p=0.030$, $r_{\mathrm{rb}}=-0.55$), whereas C4 showed no statistically significant difference after correction (adjusted $p=1.000$); the pronounced C3--C4 asymmetry in absolute $F_1$ therefore persisted regardless of whether each electrode's own subset-versus-montage change was itself significant. Several parietal and occipital electrodes also reached significance despite very low absolute $F_1$ (CP2, CP5, CP6, PO3 and PO4, all adjusted $p<0.05$), while most others did not; the central electrode FC5 illustrates the correction directly, with a raw $p=0.002$ that no longer survived Bonferroni correction (adjusted $p=0.058$).

On Cao2018, 16 of the 18 electrodes tested remained significantly different after correction, and again every significant difference was negative. The only electrodes not reaching significance were, as on Internal, the frontopolar pair: Fp1 changed from $F_1=79.68\%$ to $78.80\%$ (adjusted $p=1.000$, $r_{\mathrm{rb}}=+0.05$), and Fp2 changed from $78.45\%$ to $77.05\%$ (adjusted $p=1.000$, $r_{\mathrm{rb}}=-0.18$). The remaining two of the four highest-performing electrodes, F3 and F4, did reach significance, with reductions of $-3.55$ and $-3.35$ percentage points (adjusted $p=0.009$ and $p=0.002$; $r_{\mathrm{rb}}=-0.54$ and $-0.60$). Every central, parietal and occipital electrode tested on Cao2018 reached significance (all adjusted $p<0.001$), including C3 and C4, whose reductions remained close together ($-3.42$ and $-3.56$ percentage points; $r_{\mathrm{rb}}=-0.89$ for both), consistent with the symmetry already noted between these two electrodes on this corpus. The largest absolute change of any electrode in either dataset occurred outside the frontal region, at the central electrode FT7 ($-4.49$ percentage points, adjusted $p<0.001$, $r_{\mathrm{rb}}=-0.82$).

Taken together, the electrode-level tests show that a visibly lower descriptive mean did not by itself imply a statistically significant change. Fp1 and Fp2, the highest-performing electrode in each dataset, showed no detectable paired difference after correction in either corpus, while several weaker frontal and non-frontal electrodes changed significantly, generally by a few percentage points, even where the rank-biserial correlation indicated a large effect. These tests probe a configuration effect, whether an electrode's own detection performance changes when its region operates alone rather than as part of the complete montage. That question is distinct from the spatial pattern described above, in which frontal electrodes simply detect blinks better than posterior electrodes regardless of which subset scored them.

The electrode-level tests above establish which individual electrodes were affected by operating their region independently; the region-mean analysis that follows instead aggregates those same electrode results to ask whether the region as a whole changed relative to the same region within the full montage. Electrode-level results were aggregated into a single region-mean $F_1$ per subset, following the aggregation used for Table~\ref{tab:region_performance} and Figure~\ref{fig:region_performance}: each electrode's own macro $F_1$ was first averaged over sessions, and those per-electrode means were then averaged across the electrodes the subset contained, applied here to the subset's own self-contained rerun rather than to the full-montage run. Each subset's region-mean $F_1$ was compared against the same region's electrodes scored inside the full-montage run, on the same sessions, rather than against the full-montage headline oracle $F_1$, using the same paired-testing framework as the electrode-level comparison above, a two-tailed Wilcoxon signed-rank test paired by session with the matched-pairs rank-biserial correlation as effect size, now Bonferroni-corrected over the 13 subsets compared within each dataset rather than over individual electrodes. Because Stages~A--C were rerun independently for each subset, this design measured whether isolating a region as its own detector changed how well that region's own electrodes performed; it was not a leave-one-region-out analysis and did not establish that any region was necessary to the complete system.

The region-mean comparisons showed that every anatomical subset, in both datasets, scored significantly lower than the same electrodes' full-montage region-mean $F_1$ after Bonferroni correction, with no subset in either corpus failing to reach significance (Table~\ref{tab:exp1_subset_summary}; Figure~\ref{fig:exp1_region_boxplot}). The frontal subset was the strongest-performing region-mean comparison in both corpora. On Internal, the frontal subset reached $F_1=66.92\%$ against $F_1=69.68\%$ for the same six electrodes scored inside the full montage, a paired difference of $-2.76$ percentage points ($p<0.001$, $r=-0.80$); on Cao2018, the frontal subset reached $F_1=66.61\%$ against $F_1=68.91\%$ for the same four electrodes, a paired difference of $-2.29$ percentage points ($p=0.026$, $r=-0.47$). Both differences were statistically significant, but modest in absolute terms relative to the region's own $F_1$; the Cao2018 frontal comparison additionally had the largest corrected $p$-value and the smallest effect size of any coarse region reported below.

The remaining regions showed the same shortfall pattern at much lower absolute performance levels. On Internal, the central, parietal, occipital and posterior subsets reached region-mean $F_1$ of $13.99\%$, $1.72\%$, $2.16\%$ and $2.08\%$, respectively, against $14.80\%$, $1.94\%$, $2.30\%$ and $2.18\%$ for the same electrodes in the full montage, paired differences of $-0.81$, $-0.22$, $-0.15$ and $-0.11$ percentage points (all $p<0.001$). On Cao2018, the corresponding subsets reached $34.29\%$, $15.00\%$, $6.42\%$ and $8.62\%$, against $38.13\%$, $18.70\%$, $8.24\%$ and $10.86\%$ for the full-montage electrodes, paired differences of $-3.83$, $-3.69$, $-1.83$ and $-2.24$ percentage points (all $p<0.001$).

Two conclusions follow from these region-mean comparisons. First, the spatial hierarchy established in the full-montage analysis was preserved within the self-contained subset reruns: frontal electrodes clearly outperformed the central, parietal, occipital and posterior regions in both datasets, exactly as they did within the full 32-channel run. Second, operating an anatomical region as its own self-contained detector generally reduced the performance of that region's own electrodes relative to their performance inside the full montage: every paired difference was negative in both corpora, the Internal frontal decrease ($-2.76$ points) was larger than any other Internal region's decrease, and the Cao2018 frontal decrease ($-2.29$ points) was smaller than the central ($-3.83$ points) and parietal ($-3.69$ points) decreases and close to the posterior decrease ($-2.24$ points). Matched-pairs rank-biserial correlations were substantial for every region-mean comparison ($|r|\geq0.65$) except the Cao2018 frontal subset, where the effect was smaller ($r=-0.47$); the direction of the effect, subset $F_1$ below montage $F_1$, was nonetheless consistent across every comparison in both datasets.
